\documentclass[11pt,a4paper]{article}

\usepackage[czech]{babel}
\usepackage[T1]{fontenc}
\usepackage[utf8]{inputenc}
\usepackage{lmodern}
\usepackage[a4paper,top=1.35cm,bottom=1.35cm,left=1.7cm,right=1.7cm]{geometry}
\usepackage{enumitem}
\usepackage{amsmath}

\setlength{\parindent}{0pt}
\setlength{\parskip}{3pt}
\setlist[itemize]{leftmargin=1.3em,itemsep=1pt,topsep=1pt}

\begin{document}


\begin{center}
    {\Large\textbf{Executive Summary}}\\[1pt]
    {\normalsize Návrh osevního plánu pro rok 2025}
\end{center}

\begin{center}
    \textbf{Tým B} -- Vanessa Bergová, Daniel Garba, Aneta Lipová, Petr Lupínek, Filip Stupár
\end{center}

\vspace{-0.2cm}

\section*{1. Klíčové použité metody}

Řešení propojuje predikci cen a hektarových výnosů 20 plodin s optimalizací
osevního plánu pro modelovou farmu o rozloze 1\,000 ha. Predikce vychází
z historických dat do roku 2024 a kromě bodových odhadů zachycuje také
nejistotu pomocí 5\,000 simulovaných scénářů cen a výnosů.

Pro ceny byly porovnány naivní predikce, drift, autoregresní modely,
modely s makroekonomickými proměnnými a kombinace predikcí. U výnosů byly
testovány průměry posledních let, lineární a robustní trendy, model
s meteorologickými proměnnými a jejich kombinace. Modely byly hodnoceny
pomocí rolovacího out-of-sample backtestu podle MAE a CRPS. Finálně byla
zvolena kombinace naivní predikce a autoregresního modelu u cen a kombinace
trendu s pětiletým průměrem u výnosů.

Nejistota byla generována vyhlazeným bootstrapem historických predikčních
chyb při zachování společného vývoje jednotlivých plodin. Následná
optimalizace využívá lineární model mean--CVaR, který kombinuje očekávaný
zisk s výsledky v 10\,\% nejhorších scénářů. Jako doplňková varianta bylo
ověřeno také dělení produkce do konkrétních polí s omezením velikosti pole
a započtením mezí.

\vspace{-0.2cm}

\section*{2. Nejdůležitější dosažené výsledky}

Ceny plodin se ukázaly jako obtížně predikovatelné. Nejlepší model snížil
MAE přibližně o 5\,\% oproti naivní predikci. U výnosů byla predikce
úspěšnější; kombinace trendu a průměru dosáhla nejnižší MAE. Testované
roční meteorologické proměnné predikci výnosů nezlepšily, a proto nebyly
do finálního modelu zahrnuty.

Pro nastavení rizikové averze $\lambda=0{,}5$ vychází plán tvořený zejména
250 ha řepky, 250 ha kukuřice, 226 ha ječmene a 74 ha brambor. Z maximálně
200 ha zeleniny připadá 78 ha na zelený česnek, 77 ha na zelí a 45 ha na
salát a čekanku. Očekávaný zisk činí přibližně 44{,}6 mil. Kč a ztráta
nastává v 7\,\% simulovaných scénářů.

Rozhodovací backtest za roky 2011--2024 ukázal, že tento přístup dosáhl
průměrného zisku 28{,}8 mil. Kč ročně a v testovaném období nevykázal
ztrátový rok. Výsledky zároveň ukazují, že stabilní část plánu tvoří
především řepka, ječmen, zatímco skladba zbývajících plodin je citlivější na volbu modelu výnosů.

\vspace{-0.2cm}

\section*{3. Doporučený další postup}

V rámci dalšího rozpracování projektu je vhodné prověřit citlivost výsledného osevního plánu na volbu hodnoty $\lambda$ a porovnat tak různé úrovně důrazu na očekávaný výsledek a riziko. Zároveň lze dále rozšířit analýzu o farmově specifické údaje o výnosech a nákladech a podrobnější meteorologická data, která by umožnila ověřit, zda by přesnější vstupní data vedla ke změně predikcí a výsledného plánu. Doplňkové rozšíření o dělení polí umožňuje posoudit, zda je agregovaný osevní plán realizovatelný i při zohlednění velikosti jednotlivých polí a jejich mezí.

\vspace{-0.2cm}

\section*{4. Možné chyby v řešení}

Nejistota výsledků vyplývá především z omezené délky historických dat a z
toho, že výnosy jsou založeny na údajích za celou ČR, nikoliv za konkrétní
farmu. Skutečné riziko farmy proto může být vyšší. Náklady pro rok 2025
jsou odvozeny z roku 2024 pomocí inflace a jejich vlastní nejistota není
modelována.

Citlivostní analýza ukázala významný vliv volby modelu výnosů na skladbu
zeleniny. U některých plodin mohou být výsledky ovlivněny strukturálními
změnami v datech. Historie predikčních chyb má pouze 22 let, takže scénáře nemusí dostatečně zachytit
extrémní budoucí situace. Výsledky je proto vhodné interpretovat jako
podporu rozhodování, nikoliv jako přesnou předpověď budoucího hospodářského
výsledku.

\end{document}